\section{Introduction}
\label{sec:introduction}

Quantum software has a space problem. Useful algorithms need workspace qubits,
but present and near-term machines offer far fewer high-quality qubits than a
classical programmer would treat as ordinary memory. This pressure has produced
an important body of work on uncomputation, verified ancilla disposal, and
qubit-reuse compilation \cite{rand2019reqwire,decross2023reuse,venev2024uncomputation}.
The central optimization question is straightforward:

\begin{quote}
When has a qubit finished its job, and how can we prove that it is safe to use
again?
\end{quote}

Security introduces a second question:

\begin{quote}
If a qubit is available again, who is allowed to use it, what may they couple it
to, and what observations would reveal unauthorized use?
\end{quote}

This paper develops a bridge between those questions. The bridge begins with a
small circuit observation: in coherent teleportation, two qubits end in a known
$\ket{++}$ state and can be rotated to $\ket{00}$. That is an optimization
opportunity. The same statevector reasoning, applied adversarially to an
educational ``spy hunter'' circuit, reveals a threat-model gap: an attacker who
can access the state-preparation controls can prepare a second state for itself
without touching the signal delivered to the receiver. That is a security
lesson.

The deeper point is not that every recycled ancilla is dangerous. It is that
\emph{resource analysis and information-flow analysis must meet at the same
qubit boundary}. A compiler that proves ``this wire is now $\ket{0}$'' has
established something useful. A secure compiler should additionally establish
``this wire is uncorrelated with protected data, and no unauthorized pass can
claim it.''

\subsection{Contributions}

The paper makes six contributions:

\begin{enumerate}[leftmargin=1.8em]
  \item It gives an explicit proof that a coherent teleportation circuit with
  two final Hadamards maps $\ket{\psi00}$ to $\ket{00\psi}$ and preserves
  entanglement with an arbitrary reference system.
  \item It reconstructs three supplied Quirk exports in LaTeX and analyzes their
  complete five-qubit amplitude data rather than relying on screenshots.
  \item It proves that the advanced export is equivalent to the baseline at the
  histogram level but not at the reduced-state level, then places this result in
  a four-level hierarchy of ``same output.''
  \item It resolves the semantic routing and fixed-input fifth-wire behavior in
  a branch-conditioned source-access model, identifying when the retained wire
  perfectly reveals Alice's value and when it becomes maximally mixed.
  \item It proves a simple information-disturbance boundary that explains when
  hidden ancillas can obtain information without changing a victim state.
  \item It releases a reproducible Python artifact with command-line analysis,
  automated numerical validation, multi-version continuous integration, and
  deterministic output regression checks.
\end{enumerate}

\subsection{Claim and evidence map}

The paper is deliberately explicit about where each statement comes from. This
matters because quantum-security arguments can sound convincing long before the
threat model or measurement boundary has been nailed down.

\begin{table}[ht]
\centering
\caption{Status of the paper's central claims.}
\label{tab:evidence-map}
\begin{tabularx}{\textwidth}{@{}p{0.20\textwidth}X p{0.23\textwidth}@{}}
\toprule
Status & Claim & Evidence \\
\midrule
\evidencebadge{BaselineBlue}{PROVED} & Coherent teleportation plus two final
Hadamards maps $\ket{\psi00}$ to $\ket{00\psi}$. & Symbolic derivation and
reference-system extension. \\
\evidencebadge{AdvancedTeal}{COMPUTED} & The advanced export matches the
baseline $Z$-basis histogram but not the reduced density matrix. & Exact
partial trace of the supplied 32-amplitude arrays. \\
\evidencebadge{AdvancedTeal}{COMPUTED} & In the reconstructed fixed-input model,
$q_4$ stores the spy measurement branch and reveals $v$ exactly when $e=b$. &
Exact branch enumeration, reduced states, and trace-distance tests. \\
\evidencebadge{InterceptOrange}{RECONSTRUCTED} & The circuit is best read as a
source-access or privileged-compilation example. & Machine-readable gate
columns plus author-supplied semantic interpretation. \\
\evidencebadge{ProposalPurple}{PROPOSED} & Reclaimed qubits could become hidden
workspace for a malicious compilation or control layer. & Threat model and
research agenda; not yet a hardware demonstration. \\
\bottomrule
\end{tabularx}
\end{table}

\begin{cautionbox}[The claim is intentionally narrow]
A high-capability adversary, including a malicious insider, compromised vendor,
or state actor, is a plausible threat actor. The mathematics is actor-agnostic,
however. The paper analyzes what a privileged circuit modification could do; it
does not claim evidence that any government or quantum provider has performed
such an attack.
\end{cautionbox}
